\section{Principio de Mach, propagación bidireccional y respuesta gravitatoria colectiva}
\label{chap:mach-onda-bidireccional-rev10}
\index{Mach, principio de}
\index{onda!avanzada y retardada}
\index{Doppler!compresión y expansión}
\index{gravitación!respuesta colectiva}

\subsection{Dependencia global de la respuesta local}

Sea \(\mathcal X_\infty\) el espacio de secciones compatibles, sea
\(\operatorname{res}_v:\mathcal X_\infty\to X_v\) la restricción a un
entorno local y sea \(b:\mathcal X_\infty\to B_\infty\) la traza global de
frontera. Una respuesta inercial local \(I_v:\mathcal X_\infty\to V_v\)
es \emph{machiana} cuando existe
\(\overline I_v:\operatorname{im}(b)\to V_v\) tal que
\begin{equation}
 I_v=\overline I_v\circ b,
 \label{eq:factorizacion-mach-rev10}
\end{equation}
pero \(I_v\) no factoriza sólo a través de
\(\operatorname{res}_v\). En particular, pueden existir \(x,x'\) con
\[
 \operatorname{res}_v(x)=\operatorname{res}_v(x'),\qquad
 b(x)\ne b(x'),\qquad I_v(x)\ne I_v(x').
\]
La definición formaliza el principio de Mach: la respuesta local depende de
la clase global de cierre y no de un espacio absoluto añadido
\cite{MachMechanics}. Es el mismo esquema local--global que aparece en los
cilindros de número, en la holonomía y en el transporte de memoria.

\begin{theorem}[Criterio de factorización de la respuesta por la totalidad]
\label{thm:criterio-factorizacion-machiana-rev2}
Existe un único mapa
\(\overline I_v:\operatorname{im}(b)\to V_v\) que satisface
\(I_v=\overline I_v\circ b\) si y sólo si \(I_v\) es constante sobre cada
fibra de \(b\):
\[
 b(x)=b(y)\quad\Longrightarrow\quad I_v(x)=I_v(y).
\]
Si \(b\) es sobreyectiva, la respuesta inducida queda definida de modo único
sobre todo \(B_\infty\).
\end{theorem}

\begin{proof}
La necesidad resulta de aplicar \(\overline I_v\) a la igualdad
\(b(x)=b(y)\). Para la suficiencia, dado
\(\beta\in\operatorname{im}(b)\), definimos
\(\overline I_v(\beta)=I_v(x)\) para cualquier \(x\in b^{-1}(\beta)\). La
constancia sobre las fibras hace independiente la elección del representante;
la factorización fija después el valor de \(\overline I_v\) en cada punto de
su dominio y garantiza la unicidad.
\end{proof}

La relación entre estado y coordenada queda fijada por
\[
 [x]_{\mathfrak F}\ne x,
 \qquad
 q(x)\ne x,
 \qquad
 x\in\mathcal X_\infty.
\]
El estado conserva posición, hoja, orientación, modo aditivo o
multiplicativo, cociente, acarreo, frontera, ruta y memoria; el observable
publica sólo una coordenada. El marco local correspondiente no se añade desde
fuera, sino que se induce por restricción de la totalidad compatible y de su
conexión:
\begin{equation}
 \mathfrak F_v(x)
 :=\operatorname{Res}_v\bigl(x,
          \nabla_{\rm TPK}\bigr).
 \label{eq:marco-local-totalidad-rev10}
\end{equation}
Aquí \(\mathcal X_\infty\) no designa una suma verbal de todo el universo,
sino el sistema de secciones compatibles con orientación, frontera y memoria.
La autorreferencia es coinductiva, no circular: cada sección se reconoce por
la compatibilidad de todos sus truncamientos.

Esta formulación organiza los regímenes inerciales y acelerados mediante una
única conexión endógena. Para una curva \(\gamma\), la aceleración covariante
es
\[
 a_\gamma:=\nabla_{\dot\gamma}\dot\gamma,
 \qquad
 a_\gamma=0
 \quad\Longleftrightarrow\quad
 \gamma\ \text{es una geodésica afín de }\nabla.
\]
Si \((e_a)\) es una tétrada a lo largo de \(\gamma\), la condición
\[
 \nabla_{\dot\gamma}e_a=0
\]
define un marco paralelo y no rotante respecto de la conexión. La rotación
de un marco general se mide por sus coeficientes de conexión a lo largo de
\(\gamma\), o por su desviación respecto del transporte de Fermi--Walker
cuando éste es el estándar cinemático elegido. Para un campo arbitrario
\(v\), \(D_\gamma v=0\) significa únicamente transporte paralelo; su
negación no equivale por sí sola a aceleración o rotación.

La diferencia física queda así mensurada mediante aceleración covariante,
rotación del marco, torsión, curvatura y holonomía, sin introducir un marco
privilegiado exterior. En este sentido cada sección persistente es
\emph{autoinercial}: porta el marco inducido por su propia inserción en la
totalidad HMT.

\begin{proposition}[Autoinercia, retorno y holonomía]
\label{prop:autoinercia-holonomia-rev2}
Sea \(E\to M\) el fibrado de realización, sea \(\gamma\) un lazo basado en
\(x_0\) y sea
\(\operatorname{Hol}_{\nabla}(\gamma):E_{x_0}\to E_{x_0}\) la holonomía de
la conexión inducida. Un dato \(s_0\in E_{x_0}\) transportado en paralelo
retorna exactamente si y sólo si
\[
 s_0\in\operatorname{Fix}
 \bigl(\operatorname{Hol}_{\nabla}(\gamma)\bigr).
\]
Si el estado genealógico completo se realiza en \(E_{x_0}\) mediante una
aplicación inyectiva sobre el sector comparado, esta condición equivale
también al retorno del estado completo. Sin esa inyectividad, la conclusión
se refiere exclusivamente al dato en la fibra. Si retorna la coordenada
visible y \(s_0\) queda fuera del subespacio fijo, hay retorno de fase con
cambio del estado transportado.
\end{proposition}

\begin{proof}
El transporte alrededor del lazo actúa sobre \(E_{x_0}\) mediante
\(\operatorname{Hol}_{\nabla}(\gamma)\). La igualdad entre los datos final
e inicial equivale, por definición, a que \(s_0\) sea un vector fijo del
operador. La última equivalencia se obtiene aplicando la inyectividad de la
realización al sector declarado.
\end{proof}

\subsection{\(2\) orientaciones de una misma propagación}

Para una ruta \(\gamma\) entre emisión y absorción, la inversión
\(\gamma^\vee\) intercambia los extremos y conserva la longitud de la
trayectoria. En la realización fotónica de
\cref{chap:realizaciones-nula-material}, el par
\((\gamma,\gamma^\vee)\) representa las orientaciones retardada y avanzada
de un mismo transporte nulo. Una banda de Möbius conserva esta memoria de
orientación: una vuelta intercambia las hojas y dos vueltas restituyen la
sección orientada.

Esta estructura permite formular rigurosamente la intuición de onda piloto
bidireccional. El estado material no se identifica con una onda adicional
que lo empuja desde fuera; la sección y su fase son dos lecturas del mismo
transporte enriquecido. La identificación concreta con propagadores de Green
avanzado y retardado requiere, no obstante, un entrelazador que preserve
acción, causalidad y producto interno.

\begin{proposition}[Conjugación de orientación y reversión del propagador]
\label{prop:conjugacion-propagador-bidireccional-rev2}
Sea \(H=H^*\) un operador autoadjunto, posiblemente no acotado, que satisface
\[
 HJ_E=J_EH.
\]
Sea \(C\) una involución ortogonal de la realificación cuántica que preserva
\(D(H)\) y satisface, sobre ese dominio,
\[
 CJ_EC^{-1}=-J_E,
 \qquad CHC^{-1}=H.
\]
Entonces
\[
 U(t):=\exp(-J_EHt/\hbar_{\rm int})
\]
es un grupo unitario uniparamétrico y se cumple
\[
 \boxed{CU(t)C^{-1}=U(-t)}.
\]
\end{proposition}

\begin{proof}
El operador \(A=-J_EH/\hbar_{\rm int}\) es antiadjunto y genera \(U(t)\)
por el teorema de Stone. Las hipótesis de dominio y conjugación dan
\(CAC^{-1}=-A\); el cálculo funcional, o equivalentemente la unicidad del
grupo generado, implica
\(C\exp(tA)C^{-1}=\exp(-tA)\). Si \(H\) es acotado, la misma identidad se
obtiene término a término de la serie exponencial.
\end{proof}

La involución de palabra \(C_M=-\operatorname{rev}\) proporciona el dato
genealógico de orientación. Cuando su realización satisface las hipótesis de
la proposición, las ramas partícula y antipartícula se representan mediante
propagadores de sentidos temporales conjugados: la geometría de banda fija la
orientación y el Hamiltoniano fija la dinámica.

La lectura cinemática de compresión y expansión queda fijada, para
\(|\beta|<1\), por la rapidez
\(\zeta=\operatorname{artanh}\beta\). Para una frecuencia propia
\(\omega_0\), las dos orientaciones Doppler son
\begin{equation}
 \omega_{\rm comp}=\omega_0e^{\zeta}
 =\omega_0\sqrt{\frac{1+\beta}{1-\beta}},\qquad
 \omega_{\rm exp}=\omega_0e^{-\zeta}
 =\omega_0\sqrt{\frac{1-\beta}{1+\beta}},
 \label{eq:doppler-bidireccional-rev10}
\end{equation}
y satisfacen
\[
 \omega_{\rm comp}\omega_{\rm exp}=\omega_0^2.
\]
La reciprocidad de \eqref{eq:doppler-bidireccional-rev10} es una realización
continua del patrón bicapa: una orientación comprime y la conjugada expande,
mientras su producto conserva la escala central. No se identifica por ello
el parámetro Doppler con \(C^*\), con \(R_{\rm act}\) ni con un coeficiente
de masa; los mapas entre esas coordenadas deben construirse por separado.

\subsection{Órbita elíptica, acción y \(3\) invariantes J}

La órbita de Kepler, su levantamiento de Levi--Civita y la condición de
Sommerfeld organizan la fase sobre un ciclo elíptico. En esta rama
\[
 J_\gamma=\oint_\gamma p\,\mathrm dq
\]
es la acción orbital. Este símbolo no se confunde con el carácter modular
\(J_{\rm mod}=j-744\), que pertenece a Moonshine, ni con el invariante de
Jarlskog \(J_{\rm CKM}\), que mide violación CP. Los tres objetos comparten
una propiedad abstracta ---son invariantes obtenidos tras cerrar una
composición---, pero poseen dominios, unidades y teoremas distintos.

La condición
\[
 \varepsilon_\gamma
 \exp\!\left(\frac{iJ_\gamma}{\hbar_{\rm int}}\right)=1
\]
explica por qué la holonomía de signo desplaza en media unidad la regla de
Sommerfeld. Junto con
\(\operatorname{Area}(\theta)=\hbar_{\rm HMT}^{\rm ret}\theta\), hace
visibles el cierre \(2\pi\) y la doble vuelta \(4\pi\): el primero completa
el ciclo escalar; el segundo restituye una sección de espín o de Möbius que
ha cambiado de hoja tras una sola vuelta.

\subsection{Curvatura propagante como modo colectivo de la geometría}

La variable gravitatoria primaria es el transporte: co-tétrada, conexión,
curvatura, torsión y holonomía. Una onda gravitatoria es una variación
propagante de la curvatura o de su clase de holonomía. El tipo de estado no
requiere una coordenada elemental adicional de espín dos. Un cuanto de
helicidades \(\pm2\), si aparece después de reducir las restricciones y
cuantizar las fluctuaciones, es una excitación colectiva del sector
geométrico y no un constituyente primitivo.

La lectura machiana refuerza esta organización: la perturbación local mide
un cambio de la sección global. La componente que no permanece en la
curvatura observable puede representarse como memoria local de transporte o
contorsión sólo cuando el diagrama entre TPK y Cartan conmuta; no se identifica
la memoria discreta con la torsión por una igualdad nominal. Bajo esa
condición, la onda puede propagarse sin postular un gravitón fundamental y
la respuesta local conserva la información del cierre global.

\subsection{Hipótesis métricas y condiciones de borde de la realización machiana}

La semántica temporal del TRIT asigna al retorno la memoria del pasado, al
umbral la actualización presente y a la apertura la prolongación futura. La
realización elíptica, euclídea e hiperbólica traduce esos tres regímenes en
cierre, frontera y divergencia. Esto proporciona una lectura geométrica de
la materia como persistencia enrollada y de la propagación como apertura
hiperbólica.

Interpretar una partícula cerrada como sección bidimensional colapsada, un
electrón como clausura de dos modos nulos, o un agujero negro como inversión
global de hojas son consecuencias teóricas compatibles con esa arquitectura.
No son identidades puramente combinatorias: para adquirir estatuto físico
deben realizar una métrica, resolver las ecuaciones de campo y reproducir el
horizonte, la causalidad y las condiciones de estabilidad correspondientes.
Esta separación preserva la potencia explicativa del modelo sin convertir
una interpretación geométrica en una prueba ya obtenida.

\begin{statusbox}
La factorización machiana, la monodromía de Möbius, la reciprocidad Doppler y
la condición de fase son construcciones matemáticas explícitas en sus
dominios declarados. Su identificación conjunta con ondas piloto físicas,
agujeros negros o cosmología constituye una realización teórica que debe
satisfacer los entrelazadores y ecuaciones indicados.
\end{statusbox}
